{\small\textbf{Macro environment (Zheng, 1:05) [5:30–6:35]}}\par\smallskip
What did the economy look like?
\begin{itemize}\setlength\itemsep{0pt}
\item By the third quarter of 2021, Dutch GDP was back above its pre-COVID level, and by the end of 2021 it was almost three percent larger than before the pandemic. The recovery was complete, and that is the condition DNB's framework requires before building the buffer.
\item Then came the storm. In its spring 2022 Financial Stability Report, DNB wrote that "the economic outlook has worsened due to the war in Ukraine and high inflation".
\item Energy prices hit records, inflation reached 11.6 percent in 2022, and on 21 July 2022 the ECB raised rates by fifty basis points, its first hike.
\end{itemize}
\par\smallskip
So the macro picture gave DNB arguments \textbf{both ways}.
\begin{itemize}\setlength\itemsep{0pt}
\item The recovery said: build the buffer now, while the economy can take it.
\item The storm said: wait, because higher capital requirements could hurt a weakening economy.
\end{itemize}
\par\smallskip
Higher rates also raise the burden on borrowers, and DNB warned about the debt sustainability of governments, firms and households.
